\chapter{PROJEKTOVANJE ASEMBLERSKOG JEZIKA}
\label{ch:jezik}

\napomena{Specifikacija ulaznog jezika asemblera. 6--8 strana.
Leksička pravila su već implementirana — vidjeti docstring u \texttt{src/asm/lexer.py}.
Izvor za sintaksu: \textcite{riscvasm2024}, \textcite{fsf2024}.}

\section{Sintaksa asemblerskog jezika}
\label{sec:sintaksa}
\napomena{Leksička pravila (komentari \texttt{\#}, identifikatori, brojevi dec/hex/bin,
stringovi) i gramatika u EBNF obliku. Nacrt gramatike (uskladiti sa parserom):}

\begin{lstlisting}[caption={Nacrt gramatike asemblerskog jezika (EBNF)},label={lst:ebnf}]
program     = { line } EOF ;
line        = [ label ] [ statement ] NEWLINE ;
label       = IDENT ":" ;
statement   = directive | instruction ;
directive   = DIRECTIVE [ operands ] ;
instruction = IDENT [ operands ] ;
operands    = operand { "," operand } ;
operand     = IDENT | INTEGER | STRING | INTEGER "(" IDENT ")" ;
\end{lstlisting}

\section{Direktive asemblera}
\label{sec:direktive}
\napomena{Tabela: \texttt{.text}, \texttt{.data}, \texttt{.word}, \texttt{.byte},
\texttt{.ascii}, \texttt{.asciiz}, \texttt{.space} — sintaksa i dejstvo na brojač lokacija.}

\section{Oznake (labeli) i simboli}
\label{sec:labele}
\napomena{Pravila imenovanja, definicija i upotreba, reference unaprijed,
greške: duplirana i nedefinisana labela.}

\section{Podržani tipovi instrukcija}
\label{sec:tipovi-instrukcija}
\napomena{Sintaksa operanada po formatu (R/I/S/B/U/J) i pseudo-instrukcije
(\texttt{nop}, \texttt{li}, \texttt{la}, \texttt{mv}, \texttt{j}, \texttt{ret}) i u šta se prevode.}

\section{Pravila pisanja programa}
\label{sec:pravila}
\napomena{Jedna naredba po liniji, labela i instrukcija u istoj liniji, velika/mala slova,
obavezne sekcije.}

\section{Primjeri asemblerskih programa}
\label{sec:primjeri-programa}
\napomena{2--3 kratka primjera sa objašnjenjem; ostali su u Prilogu~\ref{sec:prilog-primjeri}.}

\lstinputlisting[style=asm,caption={Primjer programa sa petljom (\texttt{examples/loop.asm})},label={lst:loop}]{../../examples/loop.asm}
